\documentclass[10pt,a4paper]{amsart}
\usepackage[margin=19mm]{geometry}
\usepackage{amsmath,amssymb,mathtools}
\usepackage[scr=rsfs,cal=euler]{mathalfa}
\usepackage{xfrac}
\usepackage[unicode,hidelinks,bookmarks=false]{hyperref}
\newcommand{\ie}{i.e.\ }
\newcommand{\cf}{cf.\ }
\newcommand{\resp}{respectively\ }
\newcommand{\Subsection}{\subsection}
\providecommand{\nobreakdash}{\nobreak\hbox{-}}
\setlength{\emergencystretch}{2em}
\multlinegap0pt
\newtheorem{theorem}{Theorem}
\title{$\sigma$-matching and interchangeable structures on truncated polynomial algebras}
\author{Kobiljon Abdurasulov, Jobir Adashev, Feruza Toshtemirova}
\date{Source: arXiv:2603.07362v4 (CC BY 4.0)}
\begin{document}
\maketitle

\noindent\emph{Source and licence.} $\sigma$-matching and interchangeable structures on truncated polynomial algebras. Version arXiv:2603.07362v4 (CC BY 4.0), distributed by arXiv under the Creative Commons Attribution 4.0 International licence, http://creativecommons.org/licenses/by/4.0/, as recorded in the arXiv licence metadata for this exact version and shown on the abstract page https://arxiv.org/abs/2603.07362v4. Full licence notice: \texttt{licenses/arxiv-2603.07362-CC-BY-4.0.txt}.\par\smallskip

\noindent\textit{Editorial scope.} Counted unit: the article's theorem 3.1 (source0.tex lines 278--287). The article's complete original proof of it is delivered with it (source0.tex lines 289--360, one \texttt{proof} environment of 72 lines). No mathematics is written, repaired or re-derived: the statement and the proof are the article's own lines, in the article's own notation, and no retained line is substituted or re-flowed.  No further source-local context is needed: every cross-reference of every retained line resolves inside the two delivered spans.  Nothing was omitted from any retained span, so the package carries no omission record.  No retained line contains a citation, so the wrapper carries no bibliography and no bibliography entry.  No retained line contains a figure environment, an image-inclusion command or drawing code, so no image is delivered and the package declares no asset dependency.  Editorial formatting only, in addition: an AMS \texttt{amsart} preamble that restates the article's own declarations and macros used by the retained lines, namely the environments \texttt{theorem} on separate counters, together with the standard packages those lines need.  The retained environments print in this document as Theorem 1 in order of appearance, whereas the article prints the same items with the numbering of its own full document.  The complete native source of the article as distributed in the arXiv e-print for this exact version is bundled unchanged as \texttt{sources/195938/source0.tex}.
	\begin{theorem}\label{3.1}
		%Let $\mathbb{C}_n [x]$ be a~truncated polynomial algebra. 
        A linear map $\varphi: \mathbb{C}_n [x] \to \mathbb{C}_n [x]$ is an automorphism if and only if it has the form
		\begin{equation}\label{Aut}
			\varphi(e_1 )=\sum_{i=1}^n A_i e_i,
			\qquad
			\varphi(e_i ) = \sum_{j=i}^n \left( \sum_{\substack{k_1 +\cdots+k_i =j}} A_{k_1} A_{k_2} \cdots A_{k_i} \right) e_j, \quad 2 \le i \le n,
		\end{equation}
		where $A_1 \ne 0.$
	\end{theorem}

	\begin{proof}
		Assume that $\varphi: \mathbb{C}_n [x] \to \mathbb{C}_n [x]$ is an automorphism.

		First, observe that $\mathbb{C}_n [x]$ is generated by $e_1$. Indeed, by induction,
		$e_i = e_1^i, \quad 1 \le i \le n.$
		Since $\varphi$ is linear, we can write
		$\varphi(e_1 ) = \sum_{i=1}^n A_i e_i.$

		We claim that $A_1 \ne 0$. Suppose $A_1 = 0$. Then $\varphi(e_1 ) \in \langle e_2, \dots, e_n \rangle = (\mathbb{C}_n [x])^2$, which implies
		$\varphi(\mathbb{C}_n [x]) \subseteq (\mathbb{C}_n [x])^2.$
		Hence, $\varphi$ is not surjective, contradicting the assumption that $\varphi$ is an automorphism. Therefore, $A_1 \ne 0$.

		Since $\varphi$ is an algebra homomorphism, we have
		$\varphi(e_i ) = \varphi(e_1^i ) = \varphi(e_1 )^i.$
		Now compute
		\[
			\varphi(e_1 )^i = \left( \sum_{k=1}^n A_k e_k \right)^i.
		\]
		Using associativity and the multiplication rule, we obtain
		\[
			\varphi(e_1 )^i =
			\sum_{k_1,\dots,k_i}
			A_{k_1} A_{k_2} \cdots A_{k_i} \, e_{k_1 +\cdots+k_i}.
		\]
		Grouping terms by $j = k_1 + \cdots + k_i$, we get
		\[
			\varphi(e_i )
			=
			\sum_{j=i}^n
			\left(
			\sum_{\substack{k_1 +\cdots+k_i =j}}
			A_{k_1} \cdots A_{k_i}
			\right)
			e_j.
		\]
		On the other hand, assume that $\varphi$ is defined by
		\[
			\varphi(e_1 ) = \sum_{i=1}^n A_i e_i, \quad A_1 \ne 0, \qquad \text{and} \qquad
			\varphi(e_i ) = \sum_{j=i}^n \left( \sum_{k_1 +\cdots+k_i =j} A_{k_1} \cdots A_{k_i} \right) e_j.
		\]
		We show that $\varphi$ is an automorphism.

		First, we verify that $\varphi$ preserves multiplication. Since $e_i = e_1^i$, we have
		$\varphi(e_i ) = \varphi(e_1 )^i.$
		Then for any $i,j,$ where $2 \le i+j \le n:$
		\begin{center}
			$\varphi(e_i\cdot e_j )
			=
			\varphi(e_{i+j})
			=
			\varphi(e_1 )^{i+j}.$
		\end{center}
		On the other hand,
		\begin{center}
			$\varphi(e_i )\cdot \varphi(e_j )
			=
			\varphi(e_1 )^i \cdot \varphi(e_1 )^j
			=
			\varphi(e_1 )^{i+j}.$
		\end{center}
		Thus,
		$\varphi(e_i\cdot e_j ) = \varphi(e_i )\cdot \varphi(e_j ),$
		so $\varphi$ is an algebra homomorphism.

		Next, we prove that $\varphi$ is invertible. The matrix of $\varphi$ with respect to the basis $\{e_1, \dots, e_n\}$ is upper triangular with diagonal entries
		\[
			A_1, A_1^2, \dots, A_1^n.
		\]
		Since $A_1 \ne 0$, all diagonal entries are nonzero, hence the matrix is invertible. Therefore, $\varphi$ is bijective. Thus, $\varphi$ is an automorphism. The proof is complete.

		Note that, the theorem follows from the fact that $\mathbb{C}_n [x]$ is monogenic, i.e., generated by a~single element $e_1$. Hence, every automorphism is uniquely determined by the image of this generator.
	\end{proof}

\end{document}
